\section{Introduction}
\label{sec:intro}

A language model's behavior depends in part on the identity or persona from
which it answers \citep{shanahan2023roleplay}. Modern AI assistants are
trained to be helpful, refuse harmful requests, and carry out tasks while
maintaining an AI-assistant self-representation; however, changing that internal
self-representation can change how they behave. How artificial
intelligence is portrayed in pretraining data can shape model alignment
\citep{tice2026alignment}, while fine-tuning models to claim consciousness can
induce broader changes in a model's behaviors and preferences
\citep{chua2026consciousness}. These effects must ultimately act through the
model's internal state at inference time. We ask whether the model's
representation of itself as human or AI is one such variable, and we study
its behavioral role directly.

Prior work alters these behaviors through training data or fine-tuning, which
changes many things about a model at once, making it difficult to identify
the specific internal change responsible for the behavioral shift. We instead
identify a specific linear representation of human-versus-AI identity and
intervene on it directly, steering the model's position along it at inference
time with no change to the weights. Isolating the human-AI
representation lets us ask three questions separately: what it measures, when
it forms, and what changes when we steer it.

We make four contributions. (1) On OLMo 2 checkpoints, the human-AI axis forms
within the first 3.7\% of pretraining, and instruction tuning rotates it
rather than creating the contrast; on Gemma 3 4B, instruction tuning makes AI
the default self-representation (Section~\ref{sec:formation}). (2)
Register dominates the axis's raw projection, yet register-matched steering
does not reproduce axis steering's behavioral effects on Gemma 3 4B,
leaving identity as the best-supported interpretation of what steering changes
(Section~\ref{sec:calibration}). (3) On three models from two families, a
user merely claiming to be an AI shifts the model's own self-representation
toward AI, without any steering (Section~\ref{sec:context}). (4) Steering
shifts self-report far more than behavior, while also producing
safety-relevant behavioral changes: human-ward steering raises willingness to
cause harm significantly on Gemma 3 4B and Qwen 2.5 7B and directionally on
Gemma 3 12B, and raises prompt-injection execution 16-fold on 4B
(Sections~\ref{sec:context},~\ref{sec:safety}).
